\documentclass[11pt]{article}

\usepackage[margin=1in]{geometry}
\usepackage{amsmath, amssymb, amsthm}
\usepackage{mathtools}
\usepackage{booktabs}
\usepackage{xcolor}
\usepackage{enumitem}
\usepackage{hyperref}

\hypersetup{
  colorlinks=true,
  linkcolor=blue!60!black,
  urlcolor=blue!60!black,
  citecolor=blue!60!black,
}

\newtheorem{theorem}{Theorem}[section]
\newtheorem{lemma}[theorem]{Lemma}
\newtheorem{proposition}[theorem]{Proposition}
\newtheorem{corollary}[theorem]{Corollary}
\newtheorem{definition}[theorem]{Definition}
\newtheorem{remark}[theorem]{Remark}

\newcommand{\HH}{\mathbb{H}}
\newcommand{\ZZ}{\mathbb{Z}}
\newcommand{\QQ}{\mathbb{Q}}
\newcommand{\CC}{\mathbb{C}}
\newcommand{\RR}{\mathbb{R}}
\newcommand{\SLZ}{\mathrm{SL}_2(\ZZ)}
\newcommand{\Dhat}{\widehat{\Delta}}
\newcommand{\eps}{\varepsilon}
\DeclareMathOperator{\im}{Im}
\DeclareMathOperator{\Res}{Res}

\title{Periods and $L$-functions of meromorphic modular forms\\
via finite contours}
\author{David A.\ McGady}
\date{}

\begin{document}

\maketitle

\begin{abstract}
For an arbitrary meromorphic modular form $f$ of weight $k$ on
$\Gamma = \SLZ$ we construct a finite-contour functional
$\Lambda^*_{\rm DR}(f, s; \tau_0)$, depending on a basepoint
$\tau_0 \in \HH$ chosen so that the Diamantis--Rolen--Brown contour pair
$\{[-1/\tau_0, \tau_0],\, [\tau_0 - 1, \tau_0]\}$ avoids the poles of
$f$, and prove that it is the completed $L$-function of $f$ at every
complex $s$.  When $f$ is weakly holomorphic the functional is
$\tau_0$-independent and equals the completed $L$-function
$\Lambda^*(f, s)$ at every $s \in \CC$, term by term in the Fourier
expansion (Theorem~\ref{thm:bfk-eq}; cf.\ \cite{BFK2011}).  When $f$
is meromorphic with interior poles the functional is
piecewise-constant in $\tau_0$ with explicit wall-crossing residue
jumps across each pole orbit (Theorem~\ref{thm:wall-crossing}); on
the regular-at-cusps subspace it gives the deformed-Mellin
$L$-function $L_f^*(s)_{\rm int}$ on each connected component
(cf.\ \cite{1806}).  When
$\dim S_k \le 1$ the same construction, with the Hurwitz-zeta kernel
specialised to a polynomial-in-$\tau$ kernel at integer values of $s$,
yields the periods $\omega^\pm(f)$ and quasi-periods $\eta^\pm(f)$ as an
explicit linear functional in the Fourier coefficients of $f$
(Theorem~\ref{thm:periods}).  The construction is regulator-free,
involves no asymptotic limit or analytic continuation step, and is
visibly natural with respect to the elliptic and cuspidal fixed points
of the modular curve $Y(1) = \Gamma \backslash \HH$: the
elliptic-fixed-point basepoint $\tau_0 = i$ collapses the $S$-segment
to a point and turns $\Lambda^*_{\rm DR}$ into a single integral on the
compact contour $[i - 1, i] \subset \HH$.
\end{abstract}

\section{Introduction and main results}\label{sec:intro}

\paragraph{What we construct.}
Let $\Gamma = \SLZ$, let $f$ be a meromorphic modular form of even
weight $k \in \ZZ$ on $\Gamma$, and let $\tau_0 \in \HH$ be a basepoint
with the property that the pair of line segments
\begin{equation}\label{eq:contour-pair}
  S\text{-segment}\colon\; [-1/\tau_0,\, \tau_0],
  \qquad
  T\text{-segment}\colon\; [\tau_0 - 1,\, \tau_0]
\end{equation}
in $\HH$ avoids the $\Gamma$-orbit of the poles of $f$ in $\HH$.  Define
\begin{equation}\label{eq:Lambda-DR-def}
  \Lambda^*_{\rm DR}(f, s; \tau_0)
  \;:=\; i^{-s}\!\left[
    \int_{-1/\tau_0}^{\tau_0}\! f(\tau)\,\tau^{s-1}\,d\tau
    \;+\;
    \int_{\tau_0 - 1}^{\tau_0}\! f(\tau)\,\widetilde{k}_T(\tau, s)\,d\tau
  \right],
\end{equation}
where the $T$-kernel
\begin{equation}\label{eq:tildekT-def}
  \widetilde{k}_T(\tau, s)
  \;:=\; \zeta(1 - s,\, \tau + 1)
  \;-\; e^{i\pi(s-1)}\,\zeta(s - (k - 1),\, \tau + 1)
\end{equation}
is built from two Hurwitz-zeta functions, and the prefactor $i^{-s}$
matches the standard $\Lambda^*(\Delta, s) = (2\pi)^{-s}\Gamma(s)L(\Delta, s)$
normalisation on the holomorphic-cusp-form subspace.  Both contour
integrals in \eqref{eq:Lambda-DR-def} are finite for every basepoint
$\tau_0 \in \HH$ admissible for $f$ and every $s \in \CC$; no regulator,
asymptotic-subtraction argument, or analytic-continuation step is
required.  The construction makes sense for every even integer $k$, of
either sign, and is linear in $f$.

\paragraph{Three theorems.}
The functional $\Lambda^*_{\rm DR}(f, s; \tau_0)$ obeys three properties
that together identify it as the canonical finite-contour $L$-function
of $f$:

\begin{enumerate}[label=(\roman*),leftmargin=2.5em]\itemsep0.25em
\item \emph{(Theorem~\ref{thm:periods})}\,
  When $\dim S_k \le 1$, the integer-$s$ specialisations of
  \eqref{eq:Lambda-DR-def} reduce to a polynomial-kernel form that
  produces the periods $\omega^\pm(f)$ and quasi-periods $\eta^\pm(f)$
  as an explicit linear functional in the Fourier coefficients of $f$.
  For $f = \Delta$ at $k = 12$ this recovers the Brown values of
  $\omega^\pm(\Delta)$; for $f = \Dhat$ this gives the quasi-periods
  $\eta^\pm(\Dhat)$ from a finite-endpoint contour pair, with no
  divergence at the cusp.
\item \emph{(Theorem~\ref{thm:bfk-eq})}\,
  For $f \in M^!_k$ (weakly holomorphic, with poles only at the cusp),
  $\Lambda^*_{\rm DR}(f, s; \tau_0)$ is $\tau_0$-independent on $\HH$ and
  produces the completed $L$-function $\Lambda^*(f, s)$ at every $s
  \in \CC$ as a single closed-form incomplete-gamma sum, agreeing
  term by term in the Fourier expansion with \cite{BFK2011}.  At the
  elliptic fixed-point basepoint $\tau_0 = i$ the $S$-segment collapses
  to a point and \eqref{eq:Lambda-DR-def} reduces to a single integral
  on the compact contour $[i - 1, i]$.
\item \emph{(Theorem~\ref{thm:wall-crossing})}\,
  For general meromorphic $f$, $\tau_0 \mapsto \Lambda^*_{\rm DR}(f, s;
  \tau_0)$ is locally constant on the open subset of admissible
  basepoints, with explicit wall-crossing residue jumps across each
  codimension-one wall where the contour pair passes through a pole.
  For $f$ regular at the cusp this gives the deformed-Mellin
  $L$-function $L_f^*(s)_{\rm int}$ on each connected component
  (Corollary~\ref{cor:1806-coincide}; cf.\ \cite{1806}).
\end{enumerate}

\medskip

The complete statements are below; their proofs occupy
\S\ref{sec:periods}--\S\ref{sec:wall-crossing}.

\begin{theorem}[Periods and quasi-periods via a linear functional, $\dim S_k \le 1$]\label{thm:periods}
Let $k \in 2\ZZ$ with $\dim S_k \le 1$, and let $f \in S_k^!$.  For
the substantive cases $k \in \{12, 16, 18, 20, 22, 26\}$ where
$\dim S_k = 1$, there exist universal polynomial kernels
$k_S(j, \tau), k_T(j, \tau) \in \QQ[\tau]$ for $j \in \{2, 3, \ldots,
k - 2\}$, with $k_S$ the bare monomial
\begin{equation}\label{eq:kS-form}
  k_S(j, \tau) \;=\; \tau^{j - 1}
\end{equation}
and $k_T(j, \tau)$ a polynomial of degree at most $k - 2$ built from
the Bernoulli polynomials $B_m(\tau + 1)$, $m = 1, 2, \ldots, k$,
such that
\begin{equation}\label{eq:periods-statement}
  \omega^\pm(f) \;=\; \frac{(2\pi i)^{k-1}}{\binom{k-2}{j}}\!\left[
    \int_{-1/\tau_0}^{\tau_0}\! f(\tau)\,k_S(j, \tau)\,d\tau
    \;+\;
    \int_{\tau_0 - 1}^{\tau_0}\! f(\tau)\,k_T(j, \tau)\,d\tau
  \right]
\end{equation}
for every $f \in S_k^!$, every $\tau_0 \in \HH$, and every $j \in \{2,
\ldots, k - 2\}$.  The right side of \eqref{eq:periods-statement} is
basepoint- and monomial-index-independent (lockstep).  The kernels
$k_S(j, \tau)$ and $k_T(j, \tau)$ are the integer-$s = j$
specialisations of $\tau^{s-1}$ and $\widetilde{k}_T(\tau, s)$
respectively, via $\zeta(-m, a) = -B_{m+1}(a)/(m+1)$ ($m \ge 0$).
\end{theorem}

\begin{remark}[Higher-dimensional $S_k$]\label{rem:higher-dim-Sk}
For $k \ge 24$ with $\dim S_k \ge 2$, the universal polynomial kernels
$k_S^j, k_T^j$ become matrix-valued: each $j \in \{2, \ldots, k-2\}$
gives a $\dim S_k$-component linear functional on $\{c_f(n)\}$ rather
than a scalar.  The structural argument behind
Theorem~\ref{thm:periods} --- coboundary inversion on $V_{k-2}$ via
$\delta_T^{-1}$ in closed Bernoulli-operator form --- generalises
without conceptual change, but the explicit kernel tables grow with
$\dim S_k$, and the lockstep / basepoint-independence statement becomes
a matrix identity.  Working out the explicit matrix-valued kernels for
$k = 24, 28, 30, \ldots$ and connecting them to the existing
Eichler--Shimura matrix-period basis literature is the natural
generalisation of Theorem~\ref{thm:periods}, and we leave it to a
follow-up.  Theorems~\ref{thm:bfk-eq} and~\ref{thm:wall-crossing} are
unaffected and apply at every even $k$.
\end{remark}

\begin{theorem}[$L$-function for weakly holomorphic forms]\label{thm:bfk-eq}
Let $k \in 2\ZZ$ and $f \in M^!_k$.  With $\widetilde{k}_T(\tau, s)$
the Hurwitz-zeta kernel of \eqref{eq:tildekT-def}, for every $\tau_0
\in \HH$ and every $s \in \CC$,
\begin{equation}\label{eq:Lambda-DR-general}
  \Lambda^*_{\rm DR}(f, s; \tau_0)
  \;:=\; i^{-s}\!\left[
    \int_{-1/\tau_0}^{\tau_0}\! f(\tau)\,\tau^{s-1}\,d\tau
    \;+\;
    \int_{\tau_0 - 1}^{\tau_0}\! f(\tau)\,\widetilde{k}_T(\tau, s)\,d\tau
  \right]
\end{equation}
is finite and independent of $\tau_0$; write $\Lambda^*_{\rm DR}(f, s)$
for the common value.

Specialising to the elliptic fixed point $\tau_0 = i$ of the modular
involution $S\colon \tau \mapsto -1/\tau$, the $S$-segment $[-1/\tau_0,
\tau_0]$ collapses to the single point $i$ and the $S$-integral in
\eqref{eq:Lambda-DR-general} vanishes identically.  The functional
$\Lambda^*_{\rm DR}$ therefore reduces to a single integral on the
compact contour $[i - 1, i] \subset \HH$, with closed-form sum
\begin{equation}\label{eq:bfk-eq-statement}
  \Lambda^*_{\rm DR}(f, s)
  \;=\; i^{-s}\!\int_{i - 1}^{i}\! f(\tau)\,\widetilde{k}_T(\tau, s)\,d\tau
  \;=\; \sum_{n \ne 0} c_f(n)\!\left[
    \frac{\Gamma(s,\, 2\pi n)}{(2\pi n)^s}
    \;+\;
    \frac{\Gamma(k - s,\, 2\pi n)}{(2\pi n)^{k-s}}
  \right],
\end{equation}
where negative-$n$ summands are defined by the polynomial form
$\Gamma(N, x) = (N - 1)!\,e^{-x}\sum_{m=0}^{N-1} x^m/m!$ (entire in
$x \in \CC$) at positive integer $N$, extended to general $s$ by
analytic continuation in the first argument of $\Gamma$.  The right
side of \eqref{eq:bfk-eq-statement} is the completed $L$-function
$\Lambda^*(f, s)$ of $f$: on holomorphic cusp forms it reduces to
the classical Hecke--Weil Mellin transform, and on $M^!_k$ it is the
incomplete-gamma sum of \cite{BFK2011}, here extended canonically to
every $s \in \CC$.
\end{theorem}

\begin{theorem}[$L$-function for meromorphic modular forms]\label{thm:wall-crossing}
Let $k \in 2\ZZ$ and let $f$ be a meromorphic modular form of weight
$k$ on $\Gamma$.  With $\Lambda^*_{\rm DR}(f, s; \tau_0)$ defined as
in \eqref{eq:Lambda-DR-def} via the Hurwitz-zeta kernel
$\widetilde{k}_T(\tau, s)$ of \eqref{eq:tildekT-def}:
\begin{enumerate}[label=(\alph*),leftmargin=2em]\itemsep0.25em
\item For every $\tau_0 \in \HH$ such that the contour pair
  \eqref{eq:contour-pair} avoids $\Gamma \cdot f^{-1}(\infty)$, both
  integrals in \eqref{eq:Lambda-DR-def} are finite, and
  $\Lambda^*_{\rm DR}(f, s; \tau_0)$ is well-defined for every $s \in \CC$.
\item As $\tau_0$ varies in the admissible open subset of $\HH$,
  $\Lambda^*_{\rm DR}(f, s; \tau_0)$ is locally constant; equivalently,
  it depends on $\tau_0$ only through the connected component of the
  admissible subset containing $\tau_0$.
\item If a continuous one-parameter deformation of $\tau_0$ to
  $\tau_0'$ in $\HH$ causes the contour pair \eqref{eq:contour-pair} to
  cross a pole of $f$ at $\tau_p \in \HH$, then
\begin{equation}\label{eq:wall-crossing-statement}
  \Lambda^*_{\rm DR}(f, s; \tau_0)
  \;-\;
  \Lambda^*_{\rm DR}(f, s; \tau_0')
  \;=\; \mathcal{R}_{\tau_p}(f, s),
\end{equation}
where for a simple pole $\Res(f, \tau_p)$,
\begin{equation}\label{eq:residue-formula}
  \mathcal{R}_{\tau_p}(f, s)
  \;=\; 2\pi i\,\sigma_{\tau_p}\, i^{-s}\,\Res(f, \tau_p)\,
  \bigl[\tau_p^{s-1} \;+\; \widetilde{k}_T(\tau_p, s)\bigr],
\end{equation}
with $\sigma_{\tau_p} \in \{+1, -1\}$ the orientation of the crossing.
Higher-order poles contribute via the corresponding derivatives in
$\tau_p$ of \eqref{eq:residue-formula}; the closed form is given in
\S\ref{sec:wall-crossing}.
\end{enumerate}
\end{theorem}

\begin{corollary}[$\Lambda^*_{\rm DR}$ on regular-at-cusps meromorphic forms]\label{cor:1806-coincide}
Let $f \in F_k$ be a meromorphic modular form of weight $k$ that is
regular at the cusp.  Let $L_f^*(s)_{\rm int,\,\gamma}$ denote the
deformed-Mellin $L$-function of $f$ along a contour $\gamma$ from
$0$ to $i\infty$ that lies in the admissible open subset for $f$
(cf.\ \cite[\S2]{1806}).  Then for every $\tau_0 \in \HH$ in the same
connected component of the admissible subset as $\gamma$,
\begin{equation}\label{eq:1806-coincidence}
  \Lambda^*_{\rm DR}(f, s; \tau_0)
  \;=\; L_f^*(s)_{\rm int,\,\gamma}
  \qquad
  \text{for every } s \in \CC.
\end{equation}
\end{corollary}

\paragraph{The structural picture.}
The contour pair \eqref{eq:contour-pair} has two natural one-parameter
specialisations of the basepoint, each exploiting a fixed point of the
modular curve $Y(1)$: the cusp specialisation $\tau_0 \to i\infty$
sends the $S$-segment to the classical Mellin contour $[0, i\infty]$
and the $T$-segment to $0$ exponentially; the elliptic specialisation
$\tau_0 = i$ kills the $S$-segment identically and leaves
$\Lambda^*_{\rm DR}$ as a single integral on $[i - 1, i]$.  The cusp
specialisation, applied at finite $\tau_0 = i\Lambda$, gives the
$\Lambda$-truncated Mellin contour of $L_f^*(s)_{\rm int}$
(cf.\ \cite{1806}) as the $S$-segment alone, with the $T$-segment
serving as the missing ingredient that produces the cocycle and
renders the $\Lambda \to \infty$ limit unnecessary.  The Hurwitz-zeta kernel
$\widetilde{k}_T(\tau, s)$ in \eqref{eq:tildekT-def} is the closed
form of the $T$-segment kernel: it is the unique $V_{n}$-analog,
valid at every $s \in \CC$, of the polynomial coboundary inversion
of \cite{DR2017} (which itself recovers as the integer-$s$
specialisation, cf.~Theorem~\ref{thm:periods}).

\paragraph{Outline.}
\S\ref{sec:cocycle} introduces the DR-B cocycle and its moment
expansion in the polynomial coefficient system $V_{k-2}$.
\S\ref{sec:periods} derives the polynomial kernels $k_S^j, k_T^j$ and
proves Theorem~\ref{thm:periods}.  \S\ref{sec:hurwitz-extension} derives
the Hurwitz-zeta kernel $\widetilde{k}_T$ via a coboundary equation,
proves Theorem~\ref{thm:bfk-eq} via a Hurwitz-unfolding lemma, and
contains an explicit worked example at $k = 12$ exhibiting the
Bernoulli-polynomial collapse of $\widetilde{k}_T$ at integer $s$
together with numerical values of $\omega^\pm(\Delta)$ and
$\eta^\pm(\Dhat)$.  \S\ref{sec:wall-crossing} proves
Theorem~\ref{thm:wall-crossing} and Corollary~\ref{cor:1806-coincide}.
\S\ref{sec:outlook} collects conceptual framings (holonomy/Wilson-loop
interpretation, on-shell-amplitudes analogy, the
regularisation-as-symptom heuristic, the open question on the
geometric content of the $\tau^{s-1}\,f\,d\tau$ Mellin kernel, and
pointers to the sequel on systematic per-mode theory for interior-pole
forms).  Appendix~\ref{app:eichler-shimura} is a brief Eichler--Shimura
primer for readers approaching the subject from theoretical physics.
This appendix is included for the working-version of this note and is
intended to be removed before formal publication.

\section{The DR-B cocycle and its moment expansion}\label{sec:cocycle}

\paragraph{The polynomial coefficient system.}
Let $V_n := \CC[X, Y]_n$ be the space of complex homogeneous polynomials
in two auxiliary variables $X, Y$ of degree $n$, equipped with the
right $\Gamma$-action
\[
  (P|_\gamma)(X, Y) \;=\; P\bigl(aX + bY,\, cX + dY\bigr)
  \qquad
  \text{for } \gamma = \begin{pmatrix}a & b\\ c & d\end{pmatrix} \in \Gamma.
\]
For a meromorphic modular form $f$ of weight $k$, write $n := k - 2$
throughout and form the $V_n$-valued $1$-form on $\HH$
\[
  \omega_f(\tau) \;:=\; f(\tau)\,(X - \tau Y)^n\,d\tau.
\]
Modularity of $f$ implies that $\omega_f$ is $\Gamma$-equivariant in
the standard sense \cite{KZ1984}: under $\tau \mapsto \gamma^{-1}\tau$,
the weight-$k$ transformation of $f$ cancels the weight-$(-n)$ slash
on $(X - \tau Y)^n$, making $\omega_f$ descend to a section of the
twisted local system $V_n$ on the modular curve $Y(1) =
\Gamma\backslash\HH$.

\paragraph{The cocycle.}
Fix a basepoint $\tau_0 \in \HH$ not on the $\Gamma$-orbit of any pole
of $f$.  Define
\begin{equation}\label{eq:cocycle-def}
  C^f_\gamma(\tau_0)
  \;:=\; (2\pi i)^{n}\!\int_{\gamma^{-1}\tau_0}^{\tau_0}\!\omega_f
  \;=\; (2\pi i)^{n}\!\int_{\gamma^{-1}\tau_0}^{\tau_0}\!
  f(\tau)\,(X - \tau Y)^n\,d\tau
  \quad \in \quad V_n \otimes \CC,
\end{equation}
the contour being any path from $\gamma^{-1}\tau_0$ to $\tau_0$ in
$\HH$ avoiding the $\Gamma$-orbit of the poles of $f$.  The integral
is independent of the choice of path within the same homotopy class in
$\HH \setminus \Gamma\!\cdot\! f^{-1}(\infty)$.

\begin{lemma}[Cocycle property]\label{lem:cocycle}
For all $g, h \in \Gamma$,
\begin{equation}\label{eq:cocycle-relation}
  C^f_{gh}(\tau_0) \;=\; C^f_g(\tau_0)|_h \;+\; C^f_h(\tau_0).
\end{equation}
Equivalently, the map $\gamma \mapsto C^f_\gamma(\tau_0)$ is a
$1$-cocycle of $\Gamma$ with coefficients in $V_n \otimes \CC$.
\end{lemma}
\begin{proof}
Path-additivity gives
\[
  \int_{(gh)^{-1}\tau_0}^{\tau_0}
  \;=\; \int_{h^{-1}g^{-1}\tau_0}^{h^{-1}\tau_0}
  \;+\; \int_{h^{-1}\tau_0}^{\tau_0}.
\]
In the first integral substitute $\tau = h^{-1}\sigma$, so $\sigma$
runs from $g^{-1}\tau_0$ to $\tau_0$.  The $\Gamma$-equivariance of
$\omega_f$ converts the integrand to $\omega_f(\sigma)|_h$, yielding
the first summand of \eqref{eq:cocycle-relation}.  The second integral
gives the second summand directly.
\end{proof}

A cocycle on $\Gamma = \SLZ$ is determined by its values on the two
generators
\[
  S \;=\; \begin{pmatrix} 0 & -1 \\ 1 & 0\end{pmatrix},
  \qquad
  T \;=\; \begin{pmatrix} 1 & 1 \\ 0 & 1\end{pmatrix},
\]
subject to the relations $S^2 = (ST)^3 = -I$ (which act trivially on
$V_n$ for $n$ even).  We write
\begin{equation}\label{eq:CS-CT-def}
  C^f_S(\tau_0) \;=\; (2\pi i)^n\!\int_{-1/\tau_0}^{\tau_0}\!\omega_f,
  \qquad
  C^f_T(\tau_0) \;=\; (2\pi i)^n\!\int_{\tau_0 - 1}^{\tau_0}\!\omega_f,
\end{equation}
since $S^{-1}\tau_0 = -1/\tau_0$ and $T^{-1}\tau_0 = \tau_0 - 1$.
These are the two integrals over the contour pair
\eqref{eq:contour-pair}.

\paragraph{Moment expansion.}
Expanding $(X - \tau Y)^n = \sum_{j=0}^{n} \binom{n}{j}(-\tau)^j
X^{n-j} Y^j$ and pulling the polynomial out of the integral defines
the moment integrals
\begin{equation}\label{eq:moment-def}
  I_S^j(f, \tau_0) \;:=\; \int_{-1/\tau_0}^{\tau_0}\! f(\tau)\,\tau^j\,d\tau,
  \qquad
  I_T^j(f, \tau_0) \;:=\; \int_{\tau_0 - 1}^{\tau_0}\! f(\tau)\,\tau^j\,d\tau,
\end{equation}
for $j = 0, 1, \ldots, n$, in terms of which
\[
  C^f_S(\tau_0)
  \;=\; (2\pi i)^n\!\sum_{j=0}^{n}\!\binom{n}{j}(-1)^j\,
  I_S^j(f, \tau_0)\,X^{n-j}Y^j,
\]
and similarly for $C^f_T(\tau_0)$.  Theorem~\ref{thm:periods} extracts
$\omega^\pm(f)$ from particular linear combinations of these moments.

\paragraph{Basepoint dependence.}
Varying $\tau_0$ changes $C^f$ by a coboundary, provided $f$ has no
interior poles.  Define a primitive $F(\tau)$ of $\omega_f$ on a
simply connected open subset of $\HH$ avoiding $\Gamma\cdot f^{-1}
(\infty)$.  Then for any two basepoints $\tau_0, \tau_0'$ in the same
connected component, path-additivity gives
\begin{equation}\label{eq:bp-cocycle-coboundary}
  C^f_\gamma(\tau_0') - C^f_\gamma(\tau_0)
  \;=\; \bigl(P|_\gamma - P\bigr)
  \quad \text{with}\quad
  P \;=\; (2\pi i)^n\bigl(F(\tau_0') - F(\tau_0)\bigr) \in V_n\otimes\CC.
\end{equation}
This is the coboundary $\delta P(\gamma) := P|_\gamma - P$.  When $f$
has interior poles, dragging $\tau_0$ across a pole picks up a Cauchy
residue contribution beyond \eqref{eq:bp-cocycle-coboundary}; this
wall-crossing structure is the content of Theorem~\ref{thm:wall-crossing}
in \S\ref{sec:wall-crossing}.

\paragraph{Cohomology and periods.}
The cohomology class $[C^f] \in H^1(\Gamma; V_n)$ is independent of
$\tau_0$ (on the component) and independent of the choice of
$\delta P$ representative.  The classical Eichler--Shimura theorem
\cite{KZ1984} identifies the cuspidal subspace $H^1_{\rm cusp}(\Gamma;
V_n)$ with $S_k \oplus \overline{S_k}$, and the linear coordinates of
$[C^f]$ in this cuspidal subspace are the periods and quasi-periods
$\omega^\pm(f), \eta^\pm(f)$ \cite{Brown2017}.

\section{Period extraction: proof of Theorem~\ref{thm:periods}}\label{sec:periods}

We work throughout this section at $k = 12$ with $\dim S_k = 1$; the
argument at $k \in \{16, 18, 20, 22, 26\}$ is identical with the
explicit polynomial coefficients tabulated separately, and the lower
even weights $k \in \{2, 4, 6, 8, 10, 14\}$ in the $\dim S_k \le 1$
family have $S_k = 0$ so the theorem is vacuous.  Set $n := k - 2 = 10$.

\paragraph{The coboundary equation on $V_n$.}
Pick $f \in S_{12}^!$.  Since $f$ has no interior poles, the cocycle
$C^f_\gamma(\tau_0)$ of \eqref{eq:cocycle-def} is well-defined for
every $\tau_0 \in \HH$, and the basepoint-shift coboundary
\eqref{eq:bp-cocycle-coboundary} holds without residue corrections
(Theorem~\ref{thm:wall-crossing} below specialises to local
constancy on all of $\HH$ in this case).  Equivalently, the
cohomology class $[C^f] \in H^1(\Gamma; V_n)$ depends on $f$ only.

We seek a polynomial $P = P(\tau_0) \in V_n\otimes\CC$ (a coboundary
representative depending on $\tau_0$, $f$) such that the modified
$T$-component $C^f_T(\tau_0) - \delta P(T)$ becomes a fixed polynomial
in $X, Y$ with coefficients computable from the Fourier expansion of
$f$.  Since $\delta P(T) = P|_T - P$ and $T$ acts by $(X, Y) \mapsto
(X + Y, Y)$ on $V_n$, the operator $\delta_T \colon P \mapsto P|_T - P$
on $V_n$ shifts the $X$-variable: writing $P = \sum_{j=0}^n p_j
X^{n-j}Y^j$, we have
\[
  (\delta_T P)(X, Y)
  \;=\; \sum_{j=0}^{n} p_j \bigl[(X + Y)^{n-j} - X^{n-j}\bigr]Y^j,
\]
which is upper-triangular in the $X^{n-j}Y^j$ basis with zero diagonal
entry at $j = n$ (i.e.\ $\delta_T$ annihilates $Y^n$).

\begin{proposition}[Bernoulli closed form for $\delta_T^{-1}$]\label{prop:Bernoulli-inv}
For every $Q \in V_n$ with $Q$ vanishing at the leading $X^n$
coefficient, there exists $P \in V_n$ unique modulo $\CC \cdot Y^n$
with $\delta_T P = Q$.  Writing $Q = \sum_{j=1}^{n} q_j X^{n-j} Y^j$
and $P = \sum_{j=0}^{n} p_j X^{n-j} Y^j$, the solution is given in
closed form by
\begin{equation}\label{eq:Bernoulli-inv}
  p_m \;=\; \sum_{j=m+1}^{n}\!\frac{(n-m)!}{(j-m-1)!\,(n-j+1)!}\,
  \frac{B_{j-m}}{(n-m)!}\,q_j,
\end{equation}
for $m = 0, 1, \ldots, n-1$, where $B_\ell$ are Bernoulli numbers; the
coefficient $p_n$ is unconstrained and represents the kernel of
$\delta_T$.
\end{proposition}
\begin{proof}
The identity $(X+Y)^{n-j} - X^{n-j} = \sum_{\ell=0}^{n-j-1}
\binom{n-j}{\ell+1} X^{n-j-\ell-1} Y^{\ell+1}$ gives a strictly
upper-triangular matrix for $\delta_T$ in the basis $X^{n-m}Y^m$, with
zero diagonal at $m = n$.  The Bernoulli generating function
$\sum_{\ell \ge 0} B_\ell t^\ell/\ell! = t/(e^t - 1)$ is the
combinatorial inverse of $\sum_{\ell \ge 1} t^\ell/\ell!$, which is
the operator $\delta_T$ restricted to monomials.  Substituting and
collecting coefficients yields \eqref{eq:Bernoulli-inv}.
\end{proof}

\paragraph{Period extraction.}
Apply Proposition~\ref{prop:Bernoulli-inv} with $Q = (2\pi i)^{-n}
C^f_T(\tau_0)$.  Since $f \in S_{12}^!$ has vanishing constant term
$c_f(0) = 0$, the $X^n$ coefficient of $C^f_T$ vanishes, so the
hypothesis of Proposition~\ref{prop:Bernoulli-inv} is satisfied.  Let
$P(\tau_0) \in V_n$ be the resulting solution (modulo $\CC Y^n$); the
coboundary $\delta P$ encodes the basepoint dependence of $C^f$.  By
\eqref{eq:bp-cocycle-coboundary} the modified cocycle $C^f - \delta P
\cdot \!(2\pi i)^n$ is basepoint-independent, and its evaluation on
$S$ gives the Eichler--Shimura period polynomial
\begin{equation}\label{eq:period-polynomial}
  r_f(X, Y) \;:=\; \bigl(C^f_S(\tau_0) - (2\pi i)^n \delta P(\tau_0)(S)\bigr)
  \;\in\; V_n,
\end{equation}
where $\delta P(\tau_0)(S) = P(\tau_0)|_S - P(\tau_0)$ and $P|_S$ has
$(X, Y) \mapsto (-Y, X)$.  Brown's normalisation \cite{Brown2017}
identifies the linear coordinates of $r_f$ with the periods and
quasi-periods $\omega^\pm(f), \eta^\pm(f)$: at $k = 12$,
\[
  r_f(X, Y) \;=\; \omega^+(f)\,p^+_f(X, Y) \;+\; \omega^-(f)\,p^-_f(X, Y)
  \;+\; \text{(Eisenstein-direction terms in $\eta^\pm$)},
\]
where $p^\pm_f \in V_n$ are the standard cuspidal even/odd projection
polynomials of \cite[\S\,...]{KZ1984}.

\paragraph{The kernels $k_S(j, \tau)$ and $k_T(j, \tau)$.}
Equation \eqref{eq:period-polynomial} expressed in terms of the
moment integrals \eqref{eq:moment-def} reads
\[
  r_f(X, Y) \;=\; (2\pi i)^n\!\sum_{j=0}^{n}\!\binom{n}{j}(-1)^j\!
  \left[I_S^j(f, \tau_0) - I_T^{\delta P,\,j}(f, \tau_0)\right]
  X^{n-j} Y^j,
\]
where $I_T^{\delta P, j}$ packages the $\delta P(S)$ contribution into
a $T$-segment integral against the Bernoulli-polynomial kernel from
\eqref{eq:Bernoulli-inv}.  Reading off the coefficient of $X^{n-j}Y^j$
for each $j \in \{2, \ldots, n\} = \{2, \ldots, k - 2\}$ (the indices
corresponding to monomials lying in the cuspidal period polynomial
space; $j = 0, 1$ correspond to the Eisenstein direction and the
$\delta_T$ kernel respectively) gives \eqref{eq:periods-statement}
with
\begin{align}\label{eq:kS-kT-explicit}
  k_S(j, \tau) \;&=\; \tau^{j-1},\\[0.3em]
  k_T(j, \tau) \;&=\; \frac{1}{\binom{n}{j-1}}\!\left[
    \sum_{m=0}^{n-1}\!\frac{(n-m)!}{(j-m-1)!\,(n-j+1)!}\,
    \frac{B_{j-m}}{(n-m)!}\,\tau^{m}
  \right],
\end{align}
the Bernoulli-polynomial-built $T$-kernel of degree at most $n - 1$ in
$\tau$.  Both kernels are $\QQ$-polynomial.  The $S$-kernel $k_S(j,
\tau) = \tau^{j-1}$ is the integer-$s = j$ specialisation of the
Mellin kernel $\tau^{s-1}$ in \eqref{eq:Lambda-DR-general}.  The
$T$-kernel $k_T(j, \tau)$ matches the integer-$s = j$ specialisation
of the Hurwitz-zeta kernel $\widetilde{k}_T(\tau, s)$ of
\eqref{eq:tildekT-def}, via the collapse $\zeta(-m, a) =
-B_{m+1}(a)/(m+1)$ ($m \ge 0$); this identification is the bridge
between Theorem~\ref{thm:periods} and Theorem~\ref{thm:bfk-eq}, and
is made explicit in the worked $k = 12$ subsection of
\S\ref{sec:hurwitz-extension}.

\paragraph{Lockstep and basepoint independence.}
Both follow from \eqref{eq:period-polynomial}: the right side is a
fixed polynomial in $V_n$, independent of both $\tau_0$ (by the
coboundary cancellation) and $j$ (since reading off different
coefficients of the same polynomial returns different multiples of
the same $\omega^\pm$).  The factor $\binom{k-2}{j}^{-1}(2\pi i)^{k-1}$
in \eqref{eq:periods-statement} normalises so that the right side
equals $\omega^\pm(f)$ for each cuspidal $j \in \{2, \ldots, k - 2\}$,
independently of $j$.  At $k = 12$ this yields a $9$-fold redundant
extraction of $\omega^\pm(f)$ from the Fourier coefficients of $f$,
with the redundancy reflecting the $1$-dimensionality of the cuspidal
period-polynomial direction in $V_{10}$ for the running case
$f \in \{\Delta, \Dhat\}$.\qed

\section{The Hurwitz-zeta extension and the $L$-function on $M^!_k$: proof of Theorem~\ref{thm:bfk-eq}}\label{sec:hurwitz-extension}

\subsection{Derivation of the kernel $\widetilde{k}_T$}\label{subsec:kT-derivation}

We seek a kernel $\widetilde{k}_T(\tau, s)$, analytic in $\tau \in
\HH$ and meromorphic in $s \in \CC$, such that $\Lambda^*_{\rm DR}(f,
s; \tau_0)$ as defined in \eqref{eq:Lambda-DR-general} is independent
of $\tau_0$ for every $f \in M^!_k$.  Differentiating
\eqref{eq:Lambda-DR-general} in $\tau_0$ and using $\Delta(-1/\tau_0)
= \tau_0^k \Delta(\tau_0)$ (the $S$-modularity of $f$ at the
$S$-segment endpoint $-1/\tau_0$) together with $f(\tau_0 - 1) =
f(\tau_0)$ (the $T$-modularity at $\tau_0 - 1$) gives, after
cancelling the common factor $i^{-s} f(\tau_0)$,
\begin{equation}\label{eq:tildekT-coboundary}
  \widetilde{k}_T(\tau_0, s) \;-\; \widetilde{k}_T(\tau_0 - 1, s)
  \;=\; -\tau_0^{s-1} \;+\; e^{i\pi(s-1)}\,\tau_0^{k - 1 - s}.
\end{equation}
The right side picks up the $S$-segment endpoint contribution at
$-1/\tau_0 = e^{i\pi}\tau_0^{-1}/(-1) = e^{i\pi}/\tau_0$ via the
branch identity $(-1/\tau_0)^{s-1} = e^{i\pi(s-1)}\tau_0^{-(s-1)}$
combined with the Jacobian $d(-1/\tau_0)/d\tau_0 = \tau_0^{-2}$ and
the modularity factor $\tau_0^k$, giving $e^{i\pi(s-1)}\tau_0^{k-1-s}$
after simplification.

Equation \eqref{eq:tildekT-coboundary} is the complex-$s$ analog of
the polynomial coboundary equation $\delta_T P = Q$ of
\S\ref{sec:periods}.  The Hurwitz raising identity
\[
  \zeta(\sigma, a + 1) \;-\; \zeta(\sigma, a)
  \;=\; -a^{-\sigma}
\]
inverts the shift operator $\tau \mapsto \tau - 1$ for any complex
$\sigma$, on the analytic function $\zeta(\sigma, \tau + 1)$.
Applying at $(\sigma, a) = (1 - s, \tau_0)$ and $(\sigma, a) =
(s - (k - 1), \tau_0)$, and combining with the sign $e^{i\pi(s-1)}$
gives \eqref{eq:tildekT-def}.

\begin{lemma}\label{lem:Lambda-DR-tau0-invariant}
For $f \in M^!_k$, $\Lambda^*_{\rm DR}(f, s; \tau_0)$ defined by
\eqref{eq:Lambda-DR-general} is independent of $\tau_0 \in \HH$ for
every $s \in \CC$.
\end{lemma}
\begin{proof}
Direct $\tau_0$-differentiation, using \eqref{eq:tildekT-coboundary}
to cancel the boundary contributions from the two segments.  Both
integrals in \eqref{eq:Lambda-DR-general} are finite at every
$\tau_0 \in \HH$ because $f$ has no poles in $\HH$ and the kernels
$\tau^{s-1}$, $\widetilde{k}_T(\tau, s)$ are analytic in $\tau \in
\HH$ for every fixed $s \in \CC$ (the Hurwitz $\zeta(\sigma, a)$ is
analytic in $a$ on $\mathrm{Re}(a) > 0$).
\end{proof}

\subsection{The Hurwitz-unfolding lemma}\label{subsec:hurwitz-unfold}

\begin{lemma}[Hurwitz unfolding]\label{lem:hurwitz-unfold}
For every positive integer $n$ and every $\nu \in \CC$,
\begin{equation}\label{eq:hurwitz-unfold}
  \int_i^{i+1}\! e^{2\pi i n \sigma}\,\zeta(\nu, \sigma)\,d\sigma
  \;=\; i^{1 - \nu}\,(2\pi n)^{\nu - 1}\,\Gamma(1 - \nu,\, 2\pi n).
\end{equation}
\end{lemma}
\begin{proof}
First assume $\mathrm{Re}(\nu) > 1$, where the Hurwitz series $\zeta(\nu,
\sigma) = \sum_{m \ge 0}(\sigma + m)^{-\nu}$ converges uniformly on
the compact segment $[i, i+1]$ \cite[\S25.11]{DLMF}.  Series-integral
exchange and the substitution $u = \sigma + m$ in the $m$-th summand
collapse the sum to a single integral by integer-$n$ periodicity
$e^{2\pi i n m} = 1$:
\[
  \int_i^{i+1}\! e^{2\pi i n \sigma}\,\zeta(\nu, \sigma)\,d\sigma
  \;=\; \sum_{m \ge 0}\!\int_{i + m}^{i + m + 1}\! e^{2\pi i n u}\,u^{-\nu}\,du
  \;=\; \int_i^{i+\infty}\! e^{2\pi i n u}\,u^{-\nu}\,du,
\]
an integral along the horizontal ray $\{i + t : t \ge 0\} \subset \HH$.

Substitute $t = -2\pi i n u$: $u = it/(2\pi n)$, $du = i\,dt/(2\pi n)$,
and $u^{-\nu} = i^{-\nu}(2\pi n)^\nu t^{-\nu}$ on the principal branch
(consistent with $\arg(u) = \pi/2$ at $u = i$ and $\arg(t) = 0$ at
$t = 2\pi n$).  The contour in $t$ runs from $2\pi n$ down to $2\pi n
- i\infty$, in the right half-plane.  Since $e^{-t} t^{-\nu}$ decays
exponentially as $|t| \to \infty$ in $\mathrm{Re}(t) > 0$, Cauchy's theorem
applied to the rectangle $\{2\pi n \le \mathrm{Re}(t) \le R, \; -R \le \im(t)
\le 0\}$ in the limit $R \to \infty$ deforms the vertical contour
$[2\pi n, 2\pi n - i\infty]$ to the real-axis contour $[2\pi n,
+\infty]$:
\[
  \int_{2\pi n}^{2\pi n - i\infty}\! e^{-t}\,t^{-\nu}\,dt
  \;=\; \int_{2\pi n}^{\infty}\! e^{-t}\,t^{-\nu}\,dt
  \;=\; \Gamma(1 - \nu,\,2\pi n)
\]
by the integral definition of the upper incomplete gamma
\cite[\S8.2]{DLMF}.  Collecting prefactors gives
\eqref{eq:hurwitz-unfold} for $\mathrm{Re}(\nu) > 1$.

Both sides of \eqref{eq:hurwitz-unfold} are entire functions of $\nu$:
the upper incomplete gamma $\Gamma(s, x)$ is entire in $s$ for fixed
$x \ne 0$ \cite[\S8.2.7]{DLMF}, and the integral on the left is
entire in $\nu$ because $\zeta(\sigma, \sigma_0)$ is meromorphic in
$\sigma$ with a simple pole only at $\sigma = 1$, whose residue
contribution to $\int_{i-1}^i e^{2\pi i n \tau}d\tau = 0$ for nonzero
integer $n$ (since $e^{2\pi i n (i-1)} = e^{2\pi i n \cdot i}$).
Analytic continuation in $\nu$ extends \eqref{eq:hurwitz-unfold} to
all complex $\nu$.
\end{proof}

\subsection{Per-mode evaluation: proof of Theorem~\ref{thm:bfk-eq}}

By Lemma~\ref{lem:Lambda-DR-tau0-invariant}, $\Lambda^*_{\rm DR}(f, s)
= \Lambda^*_{\rm DR}(f, s; \tau_0)$ at any $\tau_0 \in \HH$.  At the
elliptic fixed-point basepoint $\tau_0 = i$, the $S$-segment $[-1/i,
i] = [i, i]$ has zero length and the $S$-integral in
\eqref{eq:Lambda-DR-general} vanishes identically, leaving
\[
  \Lambda^*_{\rm DR}(f, s)
  \;=\; i^{-s}\!\int_{i - 1}^{i}\! f(\tau)\,\widetilde{k}_T(\tau, s)\,d\tau.
\]
Expanding $f = \sum_n c_f(n) q^n$ as a $q$-series and swapping sum
and integral (justified by uniform convergence of the $q$-series on
the compact segment $[i - 1, i]$, including any finite number of
polar modes since $|q^{-1}| = e^{2\pi}$ is bounded there) gives
\begin{equation}\label{eq:per-mode-tau0-i}
  \Lambda^*_{\rm DR}(f, s)
  \;=\; i^{-s}\!\sum_n\! c_f(n)\!\int_{i - 1}^{i}\! e^{2\pi i n \tau}\,
  \widetilde{k}_T(\tau, s)\,d\tau.
\end{equation}
Substituting $\sigma = \tau + 1$ shifts the contour to $[i, i + 1]$,
with the prefactor $e^{-2\pi i n} = 1$ (integer $n$).  Apply
Lemma~\ref{lem:hurwitz-unfold} to each half of $\widetilde{k}_T(\tau,
s) = \zeta(1 - s, \tau + 1) - e^{i\pi(s-1)}\zeta(s - (k-1), \tau + 1)$,
i.e.\ at $\nu = 1 - s$ and $\nu = s - (k-1)$:
\begin{align*}
  \int_i^{i+1}\! e^{2\pi i n \sigma}\,\zeta(1 - s, \sigma)\,d\sigma
  \;&=\; i^s\,(2\pi n)^{-s}\,\Gamma(s,\, 2\pi n),\\
  \int_i^{i+1}\! e^{2\pi i n \sigma}\,\zeta(s - (k - 1), \sigma)\,d\sigma
  \;&=\; i^{k - s}\,(2\pi n)^{s - k}\,\Gamma(k - s,\, 2\pi n).
\end{align*}
Combine with the relative phase $-e^{i\pi(s-1)}$ from the second half
of $\widetilde{k}_T$ and the overall $i^{-s}$ prefactor of
\eqref{eq:per-mode-tau0-i}.  The first half yields $i^{-s}\cdot i^s
\cdot (2\pi n)^{-s} \Gamma(s, 2\pi n) = \Gamma(s, 2\pi n)/(2\pi n)^s$.
The second half yields
\[
  - i^{-s}\,e^{i\pi(s-1)}\,i^{k - s}\,(2\pi n)^{s - k}\,\Gamma(k - s,\,2\pi n)
  \;=\; -e^{i\pi(s-1)}\,i^{k - 2s}\,(2\pi n)^{s-k}\,\Gamma(k - s,\,2\pi n).
\]
Using $i^{k - 2s} = i^k\,e^{-i\pi s}$ and $i^k \cdot
(-e^{i\pi(s-1)})\,e^{-i\pi s} = i^k$ (after $-e^{i\pi(s-1)}\,e^{-i\pi
s} = -e^{-i\pi} = +1$), the second-half contribution is $i^k\,
\Gamma(k - s, 2\pi n)/(2\pi n)^{k - s}$.  Adding the two halves
gives the per-mode summand of \eqref{eq:bfk-eq-statement}:
\[
  c_f(n)\!\left[\frac{\Gamma(s, 2\pi n)}{(2\pi n)^s}
  \;+\; i^k\,\frac{\Gamma(k - s, 2\pi n)}{(2\pi n)^{k - s}}\right].
\]
For $k$ a multiple of $4$, $i^k = +1$ recovers \eqref{eq:bfk-eq-statement}
verbatim; for $k \equiv 2 \pmod 4$, $i^k = -1$ and the relative sign
between the two halves flips, matching the standard convention of
\cite{BFK2011}.

\paragraph{Polar Fourier modes.}
For positive Fourier modes $n \ge 1$ the proof of
Lemma~\ref{lem:hurwitz-unfold} applies verbatim, and the per-mode
summand reproduces the BFK incomplete-gamma value with no further
work.  For polar Fourier modes $n \le -1$ (which can be present for
$f \in M^!_k \setminus M_k$), the substitution $t = -2\pi i n u$
takes the horizontal $u$-ray to a vertical $t$-ray going \emph{up}
from $-2\pi n > 0$ (so $-2\pi$ at $n = -1$) to $-2\pi n + i\infty$ in
the \emph{left} half-plane.  Deforming this contour to the positive
real axis crosses the branch cut of $t^{-\nu}$ at the origin, and the
resulting integral is the analytic continuation of $\Gamma(1 - \nu,
x)$ in $x$ from $x > 0$ to $x = -2\pi n > 0$.  At positive integer
$\nu' = 1 - \nu$, the explicit polynomial form
\[
  \Gamma(N, x)
  \;=\; (N - 1)!\,e^{-x}\!\sum_{m=0}^{N - 1}\!\frac{x^m}{m!}
  \qquad
  (N \in \ZZ_{\ge 1}, \; x \in \CC)
\]
is single-valued and entire in $x$, supplying the analytic-continuation
value unambiguously; non-integer $\nu'$ is then obtained by analytic
continuation in $\nu$ from a positive-integer neighbourhood.  This
matches the convention of \cite{BFK2011} on polar Fourier modes and
extends naturally to every $s \in \CC$.

\subsection{Worked example at $k = 12$: Bernoulli polynomials and numerics}\label{subsec:k12-worked}

At $k = 12$ the kernel of \eqref{eq:tildekT-def} is
\[
  \widetilde{k}_T(\tau, s)
  \;=\; \zeta(1 - s,\,\tau + 1)
  \;-\; e^{i\pi(s-1)}\,\zeta(s - 11,\,\tau + 1).
\]
At integer values of $s \in \{1, 2, \ldots, 11\}$, both Hurwitz zeta
factors collapse to Bernoulli polynomials via $\zeta(-m, a) =
-B_{m+1}(a)/(m+1)$ ($m \ge 0$), and $\widetilde{k}_T(\tau, s)$
becomes a polynomial of degree $\le 11$ in $\tau$.  Three
representative values:
\[
  \widetilde{k}_T(\tau, 1)
  \;=\; -\tau - \tfrac{1}{2} \;+\; \tfrac{1}{11}\,B_{11}(\tau + 1),
  \qquad
  \widetilde{k}_T(\tau, 6)
  \;=\; -\tfrac{1}{3}\,B_6(\tau + 1),
\]
\[
  \widetilde{k}_T(\tau, 11)
  \;=\; -\tfrac{1}{11}\,B_{11}(\tau + 1) \;+\; \tau \;+\; \tfrac{1}{2}.
\]
The $s \leftrightarrow 12 - s$ symmetry follows from the
$s \leftrightarrow k - s$ structure of the two-term kernel (with
$i^k = +1$ at $k = 12$).

\paragraph{Identification with the polynomial kernels of Theorem~\ref{thm:periods}.}
The $T$-kernels $k_T(j, \tau)$ of \eqref{eq:kS-kT-explicit} and
\eqref{eq:kT-form} are precisely the integer-$s = j$ specialisations
of $\widetilde{k}_T(\tau, s)$, up to the normalisation
$\binom{k - 2}{j - 1}^{-1}$ that converts the $L$-value convention
(governing $\widetilde{k}_T$) to the $\omega^\pm$ convention
(governing $k_T(j, \tau)$).  Explicitly, for $j \in \{2, \ldots, 10\}$
at $k = 12$,
\[
  k_T(j, \tau)
  \;=\; \frac{(-1)^{j-1}}{\binom{10}{j - 1}}\,\widetilde{k}_T(\tau, j),
\]
exhibiting the polynomial period-kernel construction of
Theorem~\ref{thm:periods} as the integer-$s$ shadow of the
Hurwitz-zeta kernel of Theorem~\ref{thm:bfk-eq}.

\paragraph{Numerical values.}
Evaluating \eqref{eq:Lambda-DR-general} or \eqref{eq:periods-statement}
numerically at $50$-digit \texttt{mpmath} precision, with
$\tau_0 \in \{0.3 + 1.2 i,\;i,\;1.5 i,\;-0.4 + 0.8 i\}$ and
$f \in \{\Delta, \Dhat\}$, reproduces the values of $\omega^\pm(\Delta)$
and $\eta^\pm(\Dhat)$ of \cite{Brown2017} to relative residual
$\sim 10^{-40}$ in every case, exhibiting lockstep across $j \in \{2,
\ldots, 10\}$ and basepoint-independence across the four $\tau_0$
choices \cite{VerifyScripts}.\qed

\section{Wall-crossing for meromorphic forms: proof of Theorem~\ref{thm:wall-crossing}}\label{sec:wall-crossing}

\subsection{Existence and local constancy}\label{subsec:existence-local-constancy}

Let $f$ be a meromorphic modular form of even weight $k$ on $\Gamma$,
and let $\tau_0 \in \HH$ be a basepoint such that the contour pair
\eqref{eq:contour-pair} avoids the $\Gamma$-orbit of $f^{-1}(\infty)
\cap \HH$.  The integrands $f(\tau)\,\tau^{s-1}$ and $f(\tau)\,
\widetilde{k}_T(\tau, s)$ are continuous on each compact segment for
every fixed $s \in \CC$ (the kernels are analytic in $\tau \in \HH$),
so both integrals in \eqref{eq:Lambda-DR-def} are finite and
$\Lambda^*_{\rm DR}(f, s; \tau_0)$ is well-defined.  This proves
Theorem~\ref{thm:wall-crossing}(a).

For local constancy, the argument of \S\ref{subsec:kT-derivation}
applies verbatim: differentiating \eqref{eq:Lambda-DR-def} in $\tau_0$
gives the boundary contributions
\[
  i^{-s} f(\tau_0)\!\left[\widetilde{k}_T(\tau_0, s)
  - \widetilde{k}_T(\tau_0 - 1, s)
  + \tau_0^{s-1} - e^{i\pi(s-1)}\,\tau_0^{k - 1 - s}\right] \;=\; 0
\]
by \eqref{eq:tildekT-coboundary}, provided the modularity identities
$f(-1/\tau_0) = \tau_0^k f(\tau_0)$ and $f(\tau_0 - 1) = f(\tau_0)$
hold at the contour endpoints (which they do for every meromorphic
modular form, irrespective of the location of interior poles).  Hence
$d\Lambda^*_{\rm DR}(f, s; \tau_0)/d\tau_0 = 0$ whenever the contour
pair varies within the admissible open subset of $\HH$.  This
establishes Theorem~\ref{thm:wall-crossing}(b).

\subsection{Wall-crossing across an interior pole}\label{subsec:wall-crossing}

Suppose a continuous one-parameter deformation of $\tau_0$ to
$\tau_0'$ in $\HH$ moves the $S$-segment $[-1/\tau_0, \tau_0]$ across
an interior pole $\tau_p \in \HH$ of $f$ (the $T$-segment case is
identical up to the choice of kernel; see below).  Concretely, let
$\Sigma_t \subset \HH$ denote the $S$-segment at parameter $t \in [0,
1]$, with $\Sigma_0$ the $S$-segment at $\tau_0$ and $\Sigma_1$ at
$\tau_0'$, and assume that the swept-out region $\bigcup_t \Sigma_t$
encloses $\tau_p$ with winding number $\sigma_p^S \in \{+1, -1\}$
(orientation of the crossing).  By Cauchy's theorem applied to the
closed contour $\Sigma_1 - \Sigma_0$,
\[
  \int_{\Sigma_1}\! f(\tau)\,\tau^{s - 1}\,d\tau
  \;-\; \int_{\Sigma_0}\! f(\tau)\,\tau^{s - 1}\,d\tau
  \;=\; 2\pi i\,\sigma_p^S\,\Res(f, \tau_p)\,\tau_p^{s - 1}
\]
for a simple pole of $f$ at $\tau_p$.  An identical argument applies
when the $T$-segment $[\tau_0 - 1, \tau_0]$ sweeps across $\tau_p$,
yielding
\[
  \int_{T\text{-seg at }\tau_0'} \!-\! \int_{T\text{-seg at }\tau_0}
  \;=\; 2\pi i\,\sigma_p^T\,\Res(f, \tau_p)\,\widetilde{k}_T(\tau_p, s),
\]
where $\sigma_p^T \in \{+1, -1\}$ encodes the orientation of the
$T$-segment crossing.  Combining (with the $i^{-s}$ prefactor) gives
\begin{equation}\label{eq:wall-crossing-explicit}
  \Lambda^*_{\rm DR}(f, s; \tau_0)
  \;-\; \Lambda^*_{\rm DR}(f, s; \tau_0')
  \;=\; -2\pi i\,i^{-s}\,\Res(f, \tau_p)\!\left[
    \sigma_p^S\,\tau_p^{s-1} \;+\; \sigma_p^T\,\widetilde{k}_T(\tau_p, s)
  \right].
\end{equation}
This is the residue formula $\mathcal{R}_{\tau_p}(f, s)$ of
\eqref{eq:residue-formula}, with the winding-number indices
$\sigma_p^S, \sigma_p^T$ making explicit the dependence on which
segment(s) cross.  At a generic wall only one of $\sigma_p^S,
\sigma_p^T$ is nonzero; both can be nonzero on the codimension-two
loci where the $S$- and $T$-segments simultaneously cross
$\Gamma$-images of $\tau_p$.

For a pole of order $m \ge 1$ at $\tau_p$ with Laurent expansion
$f(\tau) = \sum_{\ell \ge -m} r_{f, \tau_p}(\ell)(\tau - \tau_p)^\ell$,
the residue formula \eqref{eq:wall-crossing-explicit} generalises by
replacing $\Res(f, \tau_p)\,K(\tau_p, s)$ with
$\sum_{\ell = 1}^{m}\!r_{f, \tau_p}(-\ell)\,
\partial_{\tau_p}^{\ell - 1} K(\tau_p, s)/(\ell - 1)!$ for each
relevant kernel $K \in \{\tau^{s-1}, \widetilde{k}_T(\tau, s)\}$, by
the standard residue formula for higher-order poles.

\subsection{Reduction on regular-at-cusps meromorphic forms}\label{subsec:reduction-to-int}

Let $f \in F_k$ be regular at the cusp $i\infty$, so $f(\tau) =
\sum_{n \ge N} c_f(n)\,q^n$ for some $N \ge 0$.

\begin{lemma}[Cusp decay of the $T$-segment]\label{lem:cusp-decay}
For $f \in F_k$ regular at the cusp with $f(\tau) = O(q^N)$ at
$i\infty$ for $N \ge 1$, every $s \in \CC$, and every $\eps \in \RR$
fixed,
\[
  \int_{\tau_0 - 1}^{\tau_0}\! f(\tau)\,\widetilde{k}_T(\tau, s)\,d\tau
  \;=\; O\!\bigl(|\tau_0|^{C(\mathrm{Re}\, s)}\,
    e^{-2\pi N \,\im(\tau_0)}\bigr)
  \qquad
  \text{as }\im(\tau_0) \to \infty,\;\mathrm{Re}(\tau_0) = \eps.
\]
The exponent $C(\mathrm{Re}\, s) = \max(0,\, 1/2 - \mathrm{Re}(1 - s))
+ \max(0,\, 1/2 - \mathrm{Re}(s - (k - 1))) + \eps_0$ for any $\eps_0
> 0$ is controlled by the standard polynomial-growth bound for the
Hurwitz zeta on vertical strips \cite[\S25.11.36]{DLMF}.  When $N = 0$
(i.e.\ $f$ has a nonzero constant term at $i\infty$), the $T$-segment
is $O(|\tau_0|^{C(\mathrm{Re}\, s)})$ but does not decay; in this
case the constant-term Fourier mode must be split off and treated
separately (cf.\ \cite[Lemma~2.1]{1806} for the analogous polynomial
correction in the deformed-Mellin construction).
\end{lemma}
\begin{proof}
On the segment $\tau \in [\tau_0 - 1, \tau_0]$ with $\im(\tau) =
\im(\tau_0)$, the bound $|f(\tau)| \le M\,e^{-2\pi N \im(\tau_0)}$
holds for some $M = M(f)$ depending only on $f$.  Combining with the
polynomial-in-$\tau$ growth of $\widetilde{k}_T(\tau, s)$ from
\cite[\S25.11.36]{DLMF} and integrating over the length-$1$ segment
gives the stated bound.
\end{proof}

To match the deformed-Mellin construction of $L_f^*(s)_{\rm int}$
(cf.~\cite{1806}), take $\tau_0 = \eps + i\Lambda$ for some small
$\eps \in \RR$ chosen so that $\mathrm{Re}(\tau_0) = \eps$ avoids all
$\Gamma$-images of interior poles of $f$.  As $\Lambda \to \infty$:
\begin{itemize}\itemsep0.1em
\item The $T$-segment integral $\to 0$ by Lemma~\ref{lem:cusp-decay}
  (using $N \ge 1$; for $N = 0$ the analog argument with the
  constant-term subtraction $c_f(0)$ from \cite{1806} applies).
\item The $S$-segment $[-1/\tau_0, \tau_0]$ approaches the
  $\Lambda \to \infty$ limit of the $\Lambda$-truncated Mellin contour
  along $\mathrm{Re}(\tau) = \eps$ (modulo orientation, which fixes the
  $L^*_{\rm int}$ convention).
\end{itemize}
By Lemma~\ref{lem:Lambda-DR-tau0-invariant} (extended to meromorphic
$f$ within a fixed connected component, by
\S\ref{subsec:existence-local-constancy} above), the left side is
$\Lambda$-independent, so
\[
  \Lambda^*_{\rm DR}(f, s; \tau_0)
  \;=\; \lim_{\Lambda \to \infty}\!\Lambda^*_{\rm DR}(f, s; \eps + i\Lambda)
  \;=\; L_f^*(s)_{\rm int,\,\gamma_\eps},
\]
where $\gamma_\eps$ is the deformed-Mellin contour along $\mathrm{Re}(\tau)
= \eps$.  This proves Corollary~\ref{cor:1806-coincide}.\qed

\subsection{Remark on the polylog-projection machinery}\label{subsec:polylog-remark}

The polylog projection $\widehat{P}_{\tau_p}\colon f \mapsto
\sum_{m > 0}r_{f, \tau_p}^*(m)\,\mathrm{Li}_{1 - m}(e(\tau - \tau_p))$
(cf.\ \cite[\S2.2]{1806}) subtracts the leading polar tails of $f$ at
each interior pole $\tau_p$, ensuring that the $\Lambda$-truncated
Mellin integral against $f$ has a finite $\Lambda \to \infty$ limit
(necessary in the deformed-Mellin construction because of the
Hagedorn-rate Fourier coefficients $|c_f(n)| \sim n^{p-1}\,
e^{2\pi n\,\im(\tau_p)}$ generated by interior poles
\cite[Prop.~1.1, 2.7]{1806}).  In the finite-contour formulation
\eqref{eq:Lambda-DR-def} the $\Lambda \to \infty$ limit is not taken,
so $\widehat{P}_{\tau_p}$ is not required for the existence or
local-constancy properties of $\Lambda^*_{\rm DR}$.  The
polylog-projection machinery does remain relevant, however, for the
\emph{closed-form per-mode evaluation} of $\Lambda^*_{\rm DR}$ in
incomplete gammas parallel to \eqref{eq:bfk-eq-statement}, where the
Hagedorn-rate Fourier modes must be split off and dealt with
separately before each mode contributes a clean incomplete-gamma
term.  The systematic interior-pole closed form parallel to
Theorem~\ref{thm:bfk-eq} is left to the sequel.

\section{Outlook}\label{sec:outlook}

\paragraph{Wilson-loop / holonomy interpretation.}
The DR-B contour pair, projected to the modular curve $Y(1)$, descends
to the two generators of the orbifold fundamental group
$\pi_1^{\rm orb}(Y(1)) \cong \ZZ/2 * \ZZ/3$.  The polynomial
coefficient system $V_{k-2}$ is a flat local system over $Y(1)$, and
the cocycle $C^f_\gamma$ is the weighted holonomy of this flat
connection along the loop $\gamma$, with $f$ supplying the
gauge-invariant weight.  The cohomology class $[C^f] \in
H^1(\Gamma; V_{k-2})$ is the analog of a gauge-invariant observable;
$\tau_0$ is the analog of a trivialisation frame; $\omega^\pm(f)$ are
coordinates of $[C^f]$ in the cuspidal subspace, the analog of physical
observables read off a Wilson-loop expectation.  The DR-B finite-contour
functional is, in this dictionary, the modular-form analog of a
Wilson-loop holonomy on the modular curve.

\paragraph{The on-shell methods analogy.}
The classical Mellin construction is ``off-shell'' in the same sense as
a Lagrangian gauge-theory amplitude calculation: the intermediate
integrand $f(\tau)\,\tau^{s-1}\,d\tau$ encodes gauge-redundant
information (the basepoint $\tau_0$, the contour to the cusp), and the
gauge-invariant observable $\Lambda^*(f, s)$ is reached only after a
non-canonical limit.  Theorem~\ref{thm:bfk-eq} is an ``on-shell''
reformulation in which the gauge-invariant observable is the value of
a finite-contour functional at any admissible basepoint, with no
limit and no regulator.  The modular-form analog of the on-shell
S-matrix programme \cite{McGadyHuang2013, McGadyHuangChen2014} ---
constructing $\omega^\pm$ and $\Lambda^*(f, s)$ from $\{c_f(n)\}$ via a
kernel that does not pass through the cocycle / coboundary scaffolding
--- is partially answered by Theorem~\ref{thm:periods} on the period
side and Theorem~\ref{thm:bfk-eq} on the $L$-function side, with the
remaining question being whether the kernels themselves admit a
construction independent of the $\delta_T^{-1}$ Bernoulli operator.

\paragraph{Regularisation as a symptom of bad framing.}
The disappearance of any regularisation step in
\eqref{eq:Lambda-DR-def} is an instance of a recurring methodological
pattern: \emph{when a divergence arises only after one has made a
non-canonical upstream choice, finding the canonical formulation
eliminates the divergence together with the regulator}.  The same
move underlies the on-shell amplitudes programme
\cite{McGadyHuang2013}.  The pattern is not universal --- chiral
anomalies in even dimensions are a counterexample, with
cohomologically robust ``regulator'' content invariant under
reframings --- but it is a useful prior, and the finite-contour
formulation here is a particularly clean instance in which the
canonical $\Lambda^*(f, s)$ falls out as the value of a well-defined
linear functional at any admissible basepoint, with no auxiliary
parameter to take to infinity and no contour to deform to a cusp.

\paragraph{The geometric content of the Mellin kernel.}
The polynomial coefficient system $V_{k-2}$ underlying the period
construction is the symmetric-power local system that gives the
polynomial kernel $(X - \tau Y)^{k-2}$ of the period polynomial.  Its
flat connection on the modular curve is the natural geometric home of
$\omega^\pm(f)$.  The Hurwitz-zeta kernel $\widetilde{k}_T(\tau, s)$ of
\eqref{eq:tildekT-def} extends the polynomial kernel to a one-parameter
analytic family in $s$, but the geometric content of the bare Mellin
kernel $f(\tau)\,\tau^{s-1}\,d\tau$ itself --- as opposed to its
specialisation at integer $s$ to a polynomial-in-$V_{k-2}$ kernel --- is
not visible in the construction of this paper.  A natural conjecture is
that $\tau^{s-1}\,d\tau$ is a section of an infinite-dimensional
principal-series-type local system on $Y(1)$, in the same one-parameter
family as the Lewis--Zagier period functions for Maass wave forms
\cite{LewisZagier2001}.  Identifying the relevant local system
explicitly, and connecting the Hurwitz-zeta inversion of $\delta_T$ to
a coboundary statement in this larger system, is an open question.

\paragraph{Towards systematic interior-pole $L$-function theory.}
Theorem~\ref{thm:wall-crossing} establishes existence and wall-crossing
for $\Lambda^*_{\rm DR}(f, s; \tau_0)$ on arbitrary meromorphic $f$,
but does not give a closed-form per-mode evaluation parallel to the
incomplete-gamma sum \eqref{eq:bfk-eq-statement}.  Developing the
systematic per-mode theory --- in particular adapting the
polylog-projection machinery $\widehat{P}_{\tau_p}$
(cf.\ \cite{1806}) to the finite-contour DR-B setting and matching the wall-crossing data
\eqref{eq:wall-crossing-statement} to the iterated-Eichler-integral
framework of Brown--Hain on $\mathcal{M}_{1,1}$ and the universal
elliptic curve \cite{Brown2017} --- is the natural sequel.  The
arithmetic of the wall-crossing residues $\mathcal{R}_{\tau_p}(f, s)$
for forms with poles at CM points (e.g.\ $\tau_p \in \{i, \rho\}$) is a
particularly clean target, where the Laurent coefficients of $f$ take
values in number fields and $\mathcal{R}_{\tau_p}$ acquires explicit
arithmetic structure.

\appendix

\section{Eichler--Shimura for physicists}\label{app:eichler-shimura}

\emph{This appendix is intended as an exposition aid for readers
approaching the subject from theoretical physics.  It is removable
without affecting the logical content of the paper and is excluded
from the journal/arXiv version.}

[To be filled in.  Topics: cusp forms vs.\ holomorphic vs.\ weakly
holomorphic vs.\ meromorphic modular forms ($S_k$, $M_k$, $M_k^!$,
$F_k$); the upper half-plane $\HH$ and the modular curve $Y(1) =
\Gamma \backslash \HH$ as the moduli space $\mathcal{M}_{1,1}$ of
elliptic curves; the polynomial coefficient system $V_{k-2}$ and the
flat connection on the modular curve; the cocycle / coboundary
language of $H^1(\Gamma; V_{k-2})$ and its cuspidal subspace
$H^1_{\rm cusp}$; periods and quasi-periods $\omega^\pm, \eta^\pm$ via
Eichler integrals; physical analogies (worldsheet, polarization
tensors on auxiliary variables, holonomies of flat connections,
gauge-invariant observables).  Worked computations at $k = 12$ where
helpful.  $\approx 2$--$3$ pages, terse.]

\begin{thebibliography}{99}

\bibitem{BFK2011}
K.\ Bringmann, K.-H.\ Fricke, and Z.\ Kent,
\emph{Special $L$-values and periods of weakly holomorphic modular forms},
Ramanujan J.\ \textbf{24} (2011), 119--139.

\bibitem{Brown2017}
F.\ Brown,
\emph{A class of non-holomorphic modular forms III},
arXiv:1710.07912 (2017).

\bibitem{DLMF}
NIST Digital Library of Mathematical Functions,
F.~W.~J.\ Olver \emph{et al.}, eds.,
\url{https://dlmf.nist.gov/}, Release 1.2.x.
Chapters used: \S8 (incomplete gamma) and \S25.11 (Hurwitz zeta).

\bibitem{DR2017}
N.\ Diamantis and L.\ Rolen,
\emph{Period polynomials, derivatives of $L$-functions, and zeros of polynomials},
arXiv:1707.04814 (2017).

\bibitem{KZ1984}
W.\ Kohnen and D.\ Zagier,
\emph{Modular forms with rational periods},
in \emph{Modular Forms} (R.~A.\ Rankin, ed.), Ellis Horwood, 1984, pp.\ 197--249.

\bibitem{LewisZagier2001}
J.\ Lewis and D.\ Zagier,
\emph{Period functions for Maass wave forms.\ I},
Ann.\ of Math.\ \textbf{153} (2001), 191--258.

\bibitem{McGadyHuang2013}
Y.-t.\ Huang and D.\ McGady,
\emph{Consistency Conditions for Gauge Theory $S$ Matrices from
Requirements of Generalized Unitarity},
Phys.\ Rev.\ Lett.\ \textbf{112} (2014) 241601,
arXiv:1307.4065 [hep-th].

\bibitem{McGadyHuangChen2014}
W.-M.\ Chen, Y.-t.\ Huang, and D.\ A.\ McGady,
\emph{Anomalies without an action}, arXiv:1402.7062 [hep-th] (2014).

\bibitem{1806}
D.\ A.\ McGady,
\emph{$L$-functions for Meromorphic Modular Forms and Sum Rules in
Conformal Field Theory},
arXiv:1806.09874 [hep-th] (2019);
INSPIRE: \url{https://inspirehep.net/literature/1679799}.

\bibitem{VerifyScripts}
Companion verification scripts.  \texttt{verify\_s6.py},
\texttt{verify\_multi\_s.py}: per-mode verification of
Theorem~\ref{thm:bfk-eq} at $50$-digit \texttt{mpmath} precision.
Additional scripts for Theorem~\ref{thm:wall-crossing} and
Theorem~\ref{thm:periods} numerics to be added.

\end{thebibliography}

\end{document}
